\section{Conclusiones}
A continuación, se presentan las conclusiones obtenidas a partir del análisis de las distintas configuraciones del circuito. En ellas se relacionan los resultados experimentales con los conceptos teóricos de corriente alterna y potencia.

\subsection{Circuito puramente resistivo}

Para esta parte del trabajo se analizó el circuito \textcolor{red}{\MakeUppercase{Poner URL llave $S_R$ cerrada}}. Debido a que se trataba de un circuito compuesto únicamente por una resistencia y una fuente de tensión alterna, se esperaba que la potencia reactiva \(Q\) fuera nula. Por lo tanto, toda la potencia entregada por la fuente debía transformarse en potencia activa \(P\), medida con el vatímetro y consumida por la resistencia.

Además, al ser un circuito puramente resistivo, la tensión y la corriente debían encontrarse en fase, por lo que el factor de potencia esperado era \(1\). Sin embargo, al calcularlo experimentalmente se obtuvo un valor de \(1{,}033\), lo que representa un error aproximado del \(3{,}2\%\). Este resultado no es físicamente posible, ya que el factor de potencia no puede ser mayor que uno. Por lo tanto, como el factor de potencia se calculó mediante la relación

\[
fp = \frac{P}{S} = \frac{P}{VI},
\]

la diferencia se atribuyó a pequeñas variaciones en las lecturas y a la precisión limitada del voltímetro, amperímetro y vatímetro.

De hecho, se pudo identificar una inconsistencia en las mediciones, ya que, al tratarse de un circuito puramente resistivo, la potencia aparente \(S = VI\) debería coincidir con la potencia activa \(P\) indicada por el vatímetro. Sin embargo, el producto entre la tensión y la corriente medidas fue

\[
S = VI = 100 \cdot 0{,}6 = 60\,\mathrm{VA},
\]

mientras que el vatímetro indicó una potencia activa de \(62\,\mathrm{W}\). Por este motivo, el valor de \(P\) resultó aproximadamente un \(3{,}2\%\) mayor que \(S\), lo cual explica que el factor de potencia calculado como \(\frac{P}{VI}\) haya dado un valor superior a la unidad.


\subsection{Circuito inductivo}

\subsection{Compensacion del factor de potencia}

\subsection{Bobina con nucleo de hierro}

\subsection{Bobina con nucleo de aluminio}
\newpage